 \documentclass[opr,utf8]{sinol}
 
 \usepackage{tikz}
 
  \title{Obliczenia w chmurze}
  \id{obl}
  \signature{kpok???}
  \author{Karol Pokorski} %% autor opracowania
  \history{10.08.2018}{JPok, przygotowanie opracowania zadania}{1.00}

\begin{document}
  \begin{tasktext}%

  \section{Treść zadania}
    Opracowujący zmienił treść zadania, aby nie pojawiało się w niej słowo "zadanie".
    Kojarzy się ono z~zadaniem, które można zakończyć (w określonym czasie),
    co błędnie może naprowadzać na problemy szeregowania zadań na procesorach.
    W tym zadaniu przeznaczamy rdzeń raz na zawsze konkretnemu zleceniu.

    Limit pamięci zmieniony na 64MB. Rozwiązanie wzorcowe zużywa bardzo mało pamięci.

  \section{Rozwiązania wzorcowe}
    W pliku \texttt{obl.cpp} jest oryginalne rozwiązanie wzorcowe dostarczone
    przez autora. W pliku \texttt{obl2.cpp} jest rozwiązanie opracowującego,
    które jest mniej efektywne -- ma gorszą stałą. Różnica polega na operowaniu
    za każdym razem na tablicy DP pełnego rozmiaru, który jest rzędu $100\,000$,
    niezależnie od danych wejściowych. Rozwiązanie autora używa tylko tyle komórek,
    ile rzeczywiście potrzeba.

  \section{Rozwiązania błędne}
    \begin{itemize}
      \item \texttt{oblb1.cpp} Wzorcówka bez sortowania po częstotliwości. Przechodzi
          tylko grupę nr $5$.
      \item \texttt{oblb2.cpp} Metoda zachłanna kupująca procesory w~kolejności
          od najtańszego w~przeliczeniu na $1$ rdzeń, a~następnie po każdym zakupie
          próbuje wykonać jak najwięcej zadań. Zadania są również wybierane zachłannie
          od najdroższych w~przeliczeniu na $1$ rdzeń. Dostaje $0$pkt.
      \item \texttt{oblb3.cpp} Rozwiązanie wzorcowe używające zbyt małej tablicy
          DP, ma ona $10\,005$ komórek. Przechodzi grupy nr $1$, $2$ oraz $3$.
      \item \texttt{oblb4.cpp} Rozwiązanie wzorcowe z tablicą DP za małą o $1$ komórkę.
          Przechodzi wszystkie grupy poza ostatnią (nr $6$).
      \item \texttt{oblb5.cpp} Wzorcówka bez long longów. Przechodzi tylko grupę nr $5$.
      \item \texttt{oblb6.cpp} Rozwiązanie zachłanne biorące procesory o najwyższej częstotliwości
        jako pierwsze i szukające lower\_bound'em zlecenia możliwego do wynonania o możliwie
        największej częstotliwości. Rozwiązanie stworzone z myślą o podzadaniu nr $5$.
        Dostaje $0$pkt.
    \end{itemize}

  \section{Testy}
    W większości są losowe, z odpowiednim doborem zakresów rdzeni, częstotliwości oraz kosztu
    -- funkcje \texttt{SimpleGenTasks} oraz \texttt{SimpleGenProcessors}.

    W grupie nr $1$ zawarte są $2$ testy skrajne $n = m = 1$, jeden gdzie się da wykonać zlecenie
    i~drugi, gdzie się nie da -- odpowiedź to $0$.

    Funkcje \texttt{GenTasks} oraz \texttt{GenProcessors} były używane do stworzenia "klonów"
    testów losowych -- w większości grup jest seria testów generowanych losowo, do każdego
    testu zaraz po nim jest generowany jego "klon" metodą zawartą w~tych magicznych funkcjach
    dostarczonych przez autora.

    W niektórych grupach używana jest funkcja \texttt{GenFreqPriceProportional}, która rozmieszcza
    zlecenia i procesory losowo względem podanej (jako parametr) funkcji $f(x)$. Założenie jest takie,
    że wszystkie zlecenia i procesory wrzucamy do jednego wora, a następnie rysujemy wykres
    wartości (cena/nagroda) względem częstotliwości -- ten wykres to właśnie funkcja $f(x)$.
    Obserwacja jest taka, że jeśli $\forall_i c_i = C_i = 1$ oraz $f$ jest rosnąca, to nie opłaca się wykonywać żadnych zleceń.
    Jeśli natomiast jest malejąca, to nie opłaca się robić nic. W testach jest kilka funkcji
    jednego i drugiego typu. Niestety powyższa obserwacja nie jest prawdą bez zalożenia o liczbie
    rdzeni, niemniej jednak gdy wygenerujemy test z większą liczbą rdzeni na zadania i procesory,
    to trzeba "zbierać" tańsze zadania, aby było nas stać na kupno droższego procesora (gdy $f$ rosnąca),
    co też wydaje się być sensownym testem. Zawarte w grupie nr $6$.

    Test \texttt{obl6c.in} łączy w jednym teście zbiór zadań i procesorów, z których żaden się nie opłaca
    oraz część, gdzie wszystko się opłaca - stworzona jest tam funkcja nieciągła w tym celu.

  \section{Do rozważenia}
    Można wyrzucić/zmodyfikować niektóre duże testy, aby rozwiązanie bez long long'ów dostawało
    rozsądniejszą liczbę punktów -- obecnie ma $18$pkt przechodząc tylko grupę nr $5$.

    Dołożyć rozwiązanie bazujące na matchingu ważonym -- sugerowane przez autora.

  \end{tasktext}
\end{document}
